% !TEX root = ./main.tex



\section{Compilation and Inference Algorithm}\label{sec:inference}
This section describes the inference algorithm for closed and well-formed \NDICE programs in several steps.
First, we describe how \NDICE expressions and programs are compiled into Boolean formulas. 
Afterwards, we demonstrate how these formulas can be used to represent the behavior of the compiled expression as an \textit{Algebraic Decision Diagram} (ADD).
We then lift this ADD to a \textit{Markov Decision Process} (MDP) and relate the semantics of the original \NDICE expression to conditional reachability properties on this MDP.

\subsection{Compilation of \NDICE Expressions to Boolean Formulas}\label{sec:nDice_comp}

	\noindent
	We will now present how the \DICE compilation procedure from \Cref{sec:prelims_Dice} can be extended to compile \NDICE expressions.
	In this extended compilation process, a \emph{compilation judgment} for a \NDICE expression $e$ has the form $\Gamma \vdash e \colon \tau \leadsto (\dot{\varphi},\, \gamma,\, t)$.  
	Here, the model formula $\dot{\varphi}$ and accepting formula $\gamma$ are unchanged compared to the $\DICE$ compilation.
	However, since nondeterministic behavior depends on the order of flip operations, an \textit{unordered} set of weights is no longer sufficient.
	We thus compile an annotated trace of flips $t$, which both tracks the order of the flips in $e$ and stores their behavior.
	To do so, $t$ is a list of pairs $f_i\colon a_i$, where $f_i$ is the Boolean variable which represents a flip and the label $a_i$ denotes the flip's behavior:
	We use $a_i = n$ if $f_i$ belongs to a nondeterministic flip, and $a_i = p$ if $f_i$ represents a probabilistic flip $\FLIP{p}$.
	
	The rules for \NDICE compilation of expressions can be found in \Cref{nDice Compilation}, and they are largely identical to the rules for \DICE~\cite{holtzen2020scaling} as seen in \Cref{fig:dice_Comp}.
	However, there are two notable changes:
	For one, there is an additional rule \texttt{nC-NFlip} to compile $\NFLIP$ statements. 
	This rule behaves fully analogous to the rule for probabilistic flips, but the boolean variable $f$ is now annotated with an $n$ to mark that it corresponds to a nondeterministic flip.
	Secondly, all rules now consider a trace of literals instead of a set of weights, which has little effect on most rules.
	However, when compiling a let-statement with the rule \texttt{nC-Let}, the trace of flips has to concatenate the trace for $e_1$ \textit{before} the trace for $e_2$ to capture the order of flip-operations.
	As for if-statements, the traces for the branches are combined by simply concatenating them, but any interleaving of the traces would be sound as long as the order of the individual sub-traces is preserved.



%
%
%\noindent%
%
%
\begin{figure}
	\begin{center}
		\begin{adjustbox}{max width=\textwidth}
			\renewcommand{\arraystretch}{3.2}
			\begin{tabular}{c c}
				
				%value 
				\inference[\ifthenelse{\boolean{showRuleNames}}{\texttt{(nC-Val)}}{}]
				{}{  \Gamma \vdash v \colon \tau \leadsto (v, \TRUE, [])}
				&
				
				%Variable
				\inference[\ifthenelse{\boolean{showRuleNames}}{\texttt{(nC-Ident)}}{}]
				{\Gamma(x) = \tau}{  \Gamma \vdash x \colon \tau \leadsto (F_\tau(x), \TRUE, [])}
				\\
				
				%flip
				\inference[\ifthenelse{\boolean{showRuleNames}}{\texttt{(nC-Flip)}}{}]
				{\text{variable $f$ fresh}}{  \Gamma \vdash \FLIP{\theta} \colon \BOOL \leadsto (f, \TRUE, [f\colon \theta])}
				&
				
				%nflip
				\inference[\ifthenelse{\boolean{showRuleNames}}{\texttt{(nC-NFlip)}}{}]
				{\text{variable $f$ fresh}}{  \Gamma \vdash \NFLIP \colon \BOOL \leadsto (f, \TRUE, [f\colon n])}
				\\
				
				%tuple construction
				\inference[\ifthenelse{\boolean{showRuleNames}}{\texttt{(nC-Tup)}}{}]
				{\Gamma(x_1) = \tau_1 \qquad \Gamma(x_2) = \tau_2}{  \Gamma \vdash (x_1, x_2) \colon \tau_1 \times \tau_2 \leadsto ((F_{\tau_1}(x_1), F_{\tau_2}(x_2)) , \TRUE, [])}
				&
				
				%observe
				\inference[\ifthenelse{\boolean{showRuleNames}}{\texttt{(nC-Obs)}}{}]
				{\Gamma \vdash \AEXP \colon \BOOL \leadsto (\varphi, \TRUE, [])}
				{  \Gamma \vdash \OBSERVE{\AEXP} \colon \BOOL \leadsto (\TRUE, \varphi, [])}
				\\
				
				%first
				\inference[\ifthenelse{\boolean{showRuleNames}}{\texttt{(nC-Fst)}}{}]
				{\Gamma(x) = \tau_1 \times \tau_2}
				{ \Gamma \vdash \FST x \colon \tau_1 \leadsto (F_{\tau_1}(x_l), \TRUE, [])}
				&
				
				%second
				\inference[\ifthenelse{\boolean{showRuleNames}}{\texttt{(nC-Snd)}}{}]
				{\Gamma(x) = \tau_1 \times \tau_2}
				{ \Gamma \vdash \SND x \colon \tau_1 \leadsto (F_{\tau_2}(x_r), \TRUE, [])}
				\\
				
				\multicolumn{2}{c}{
					%if
					\inference[\ifthenelse{\boolean{showRuleNames}}{\texttt{(nC-ITE)}}{}]
					{\Gamma \vdash \AEXP \colon \BOOL \leadsto (\varphi_g, \TRUE, []) 
						\qquad \Gamma \vdash e_T \colon \tau \leadsto (\dot{\varphi}_T, \gamma_T, t_T) 
						\qquad \Gamma \vdash e_F \colon \tau \leadsto (\dot{\varphi}_F, \gamma_F, t_F)}
					{  \Gamma \vdash \ITE{\AEXP}{e_T}{e_F} \colon \tau \leadsto \left( 
						(\varphi_g \overset{\tau}{\land} \dot{\varphi}_T) \overset{\tau}{\lor} (\overline{\varphi}_g \overset{\tau}{\land} \dot{\varphi}_F), \;
						(\varphi_g \land \gamma_T) \lor (\overline{\varphi}_g \land \gamma_F) ,\;
						t_T \,@\, t_F \right)}
				}\\
				
				\multicolumn{2}{c}{
					%let
					\inference[\ifthenelse{\boolean{showRuleNames}}{\texttt{(nC-Let)}}{}]
					{\Gamma \vdash e_1 \colon \tau_1 \leadsto (\dot{\varphi}_1, \gamma_1, t_1) \qquad 
						\Gamma \cup \{x \mapsto \tau_1 \} \vdash e_2 \colon \tau_2 \leadsto (\dot{\varphi}_2, \gamma_2, t_2)}
					{  \Gamma \vdash\LET{x = e_1}{e_2} \colon \tau_2 \leadsto \left(
						\dot{\varphi}_2[x\overset{\tau_1}{\mapsto} \dot{\varphi}_1], \;
						\gamma_1 \land \gamma_2[x\overset{\tau_1}{\mapsto} \dot{\varphi}_1],\;
						t_1 \,@\, t_2 \right)}
				}
				
			\end{tabular}
		\end{adjustbox}
	\end{center}
	\caption{\NDICE Compilation rules for expressions. These assume, without loss of generality but for simplicity, that fst, snd, and tuple construction are only ever performed with identifiers as arguments.}
	\label{nDice Compilation} 
\end{figure}%
%
%
%

\begin{example}
	As an example, consider the following expression and its compilation judgment:
	\begin{align*}
		\LET{x = \FLIP{\tfrac{2}{3}}}{\LET{y = \NFLIP}{\LET{z = \OBSERVE{x \lor y}}{(x \land y,\, y)}}}  \tag{ExComp}\label{ExComp} \\
		\{\}\vdash \text{\ref{ExComp}} \colon \BOOL \times \BOOL \leadsto \left((f_1 \land f_2,\, f_2),\; f_1 \lor f_2,\; \left[f_1\colon {\tfrac{2}{3}},\, f_2\colon n \right]\right)
	\end{align*}
	In the compilation result, $f_1$ represents the first flip, which is probabilistic, and $f_2$ represents the second flip, which is nondeterministic.
	This behavior (along with the order of the flips) is stored in the trace $\left[ f_1 \colon {\tfrac{2}{3}}, \; f_2\colon n\right]$.
	The model formula $(f_1 \land f_2,\, f_2)$ represents the output of the compiled expression, and each component of the model formula describes the behavior of the corresponding component in the output of \ref{ExComp}.
	Finally, the accepting formula $f_1 \lor f_2$ encodes the only observe statement in \ref{ExComp}, hence if $f_1 \lor f_2$ is satisfied, then all observations in \ref{ExComp} succeed.
	\hfill$\triangleleft$
\end{example}


\paragraph{Compilation of Functions and Programs}
We are not discussing the compilation rules for \NDICE programs at this point, since they are fully identical to the rules from \DICE~\cite{holtzen2020scaling}.
Furthermore, programs are still compiled to triples of the form $(\dot{\varphi}, \gamma, t)$, thus all subsequent inference steps are identical between the inference process for expressions and programs.
For completeness, rules for functions and programs are given in \ext{\Cref{app:function_comp}}{the extended version \citeExtended}.





\subsection{Algebraic Decision Diagrams}\label{sec:ndice_add}
This section presents how the previously compiled Boolean formulas can be used to build a decision diagram which compactly represents the compiled program.
Thereafter, we will enrich this decision diagram to an MDP on which we then do inference.

Consider a compilation judgment $\{\} \vdash e \colon \tau \leadsto (\dot{\varphi}, \gamma, t)$ for a closed \NDICE expression $e$. 
Our aim is now to use the compilation result $(\dot{\varphi}, \gamma, t)$ to define a function $(\dot{\varphi} \,|\, \gamma)_t\colon \BOOL^{len(t)}\rightarrow \llbracket\tau \rrbracket \, \cup \{R\}$, which mimics the behavior of $e$ for any given outcome of the flip operations in $e$. 
Here, the notation $\llbracket\tau \rrbracket$ refers to the set of all values of type $\tau$.

As a first observation, notice that if we assign either $\TRUE$ or $\FALSE$ to each flip operation in $e$,  then $e$ evaluates to a single value $v\in \llbracket\tau \rrbracket$ (assuming all observations in $e$ would succeed).
Similarly, we can assign $\TRUE$ or $\FALSE$ to each variable in $\dot{\varphi}$ to evaluate $\dot{\varphi}$ to a value $v\in \llbracket\tau \rrbracket$.  
Notably, for closed expressions each variable in $\dot{\varphi}$ and $\gamma$ corresponds to a flip statement in $e$, and hence each assignment to flips in $e$ defines an assignment to the variables in  $\dot{\varphi}$ and $\gamma$ and vice versa.
Now, given any assignment to flip operations in $e$, we thus want the function $(\dot{\varphi} \,|\, \gamma)_t$ to return the same value $v\in \llbracket \tau \rrbracket$ as $e$ if all observations in $e$ would succeed, and otherwise $(\dot{\varphi} \,|\, \gamma)_t$ should return $R$ to denote a "rejected" run.


Second of all, the function $(\dot{\varphi} \,|\, \gamma)_t$ takes a vector of Boolean values $(b_1,\dots,b_k) \colon \BOOL^k$ with $k=len(t)$ as input.
Thus, we need to define how $\gamma$ and $\dot{\varphi}$ are evaluated with a vector $(b_1,\dots,b_k)$ before we define the function $(\dot{\varphi} \,|\, \gamma)_t$.
In particular, we need to decide which variable in $\gamma$ and $\dot{\varphi}$ is substituted with which Boolean value $b_i$.
To this end, we use the trace of flips $t$ and substitute the $i$-th flip $f_i$ from $t$ with $b_i$, written as $\gamma \statesubst{f_i}{b_i}$.
The Boolean formulas resulting from these substitutions are ground, and can be evaluated to a Boolean value.
Thus, the Boolean formula $\gamma$ represents a function of type $\BOOL^k\rightarrow \BOOL$, and the tuple of Boolean formulas $\dot{\varphi}$ can be evaluated recursively as a function of type $\BOOL^k\rightarrow \tau$:
%
\begin{align*}
\gamma (b_1,\dots,b_k) &\coloneqq \gamma\statesubst{f_1}{b_1}\dots\statesubst{f_k}{b_k}\\
 \left(\dot{\varphi}_1,\; \dot{\varphi}_2\right)(b_1,\dots,b_k) &\coloneqq \left(\dot{\varphi}_1 (b_1,\dots,b_k), \;\dot{\varphi}_2(b_1,\dots,b_k)\right)
\end{align*}
We can now combine the two formulas into a function $(\dot{\varphi} \,|\, \gamma)_t \colon \BOOL^k\rightarrow \tau\, \cup \{R\}$ by mimicking the behavior of \NDICE expressions:
If $\gamma$ returns \TRUE, then all observations succeeded and we return the output of $\dot{\varphi}$.
If $\gamma$ returns \FALSE, then an observation failed and the run will be rejected, so we return the special value $R$.
Formally, we define this as:
$$ (\dot{\varphi} \,|\, \gamma)_t (b_1,\dots,b_k)\coloneqq \begin{cases}
	\dot{\varphi} (b_1,\dots,b_k), & \text{if } \gamma(b_1,\dots,b_k) = \TRUE\\
	R, & \text{if } \gamma(b_1,\dots,b_k) = \FALSE
\end{cases}$$%
Once this \enquote{guarding} operator $(\cdot \,|\, \cdot)$ is defined, the ADD for $(\dot{\varphi} \,|\, \gamma)_t$ can be constructed from decision diagrams for $\dot{\varphi}$ and $\gamma$ using generic algorithms \cite{fujita1997multi}.
Thus, we can represent $(\dot{\varphi} \,|\, \gamma)_t$ as an ADD, and if an ADD %$\mathcal{A}$ 
uses the variable order defined by $t$ and represents $(\dot{\varphi} \,|\, \gamma)_t$, then we call it an ADD for the compiled formulas $(\dot{\varphi}, \gamma, t)$.
If an ADD for $(\dot{\varphi}, \gamma, t)$ is a reduced ADD, then we call it the reduced ADD for $(\dot{\varphi}, \gamma, t)$.

\begin{remark}[Permissible Variable Orderings]
	As mentioned in \Cref{sec:prelims_ADD}, the variable ordering which is used during construction has a stark effect on the size of the ADD.
	This appears to call for experiments with different variable orderings to speed up inference, as was done for \DICE \cite{cheng2021flip}.
	However, changing the variable order amounts to changing the order of flips, which can alter the semantics of a \NDICE expression.
	In fact, variable orders have to adhere to the natural program order to be sound for inference, and the trace $t$ provides one such order. 
	\hfill $\triangleleft$
\end{remark}

\subsection{Lifting ADDs to Markov Decision Processes}\label{sec:MDP_inf}

Given a \NDICE expression $e$ which compiles to a triple $(\dot{\varphi}, \gamma, t)$, the last part of the inference process is to lift the previously constructed ADD for $(\dot{\varphi}, \gamma, t)$ to an MDP by re-introducing the behavior of the individual flip operations into the ADD.
To do so, we consider the variable stored in an inner ADD node, and determine the behavior of the corresponding flip operation by retrieving the annotated variable from the trace $t$.
The ADD node will then be modeled by an MDP state, which behaves either probabilistic or nondeterministic depending on the nature of the flip operation.
Afterwards, the reachability probabilities in this lifted MDP will align with the probabilities that the compiled expression $e$ returns a given value.

We define the \textit{unnormalized} MDP $\mathcal{M}^u = (S, s_{init}, \mathit{Act}, P, AP, L)$ lifted from an ADD $\mathcal{A} = (V, T, r, E)$ for $(\dot{\varphi}, \gamma, t)$ as follows: 
We define $S = V \cup T$ as the set of nodes of the ADD, $s_{init}$ as the root $r$ of the ADD, $\mathit{Act} = \{l, r, d\}$ and we define the transition relation as follows:
\begin{itemize}
	\item If $s \in T$ was lifted from a terminal node, then we only have the transition $s\xrightarrow{d, 1} s$
	\item If $s\in V$ was lifted from an inner node with probabilistic variable $f$ with $(f\colon \theta) \in t$, then we have the transitions $s\xrightarrow{d, \theta} \mathit{then}(s)$ and $s\xrightarrow{d, 1 - \theta} \mathit{else}(s)$
	\item If $s\in V$ was lifted from an inner node with nondeterministic variable $f$ with $(f\colon n) \in t$, then we have the transitions $s\xrightarrow{l, 1} \mathit{then}(s)$ and $s\xrightarrow{r, 1} \mathit{else}(s)$
\end{itemize}
Above, we re-use the notations $\mathit{then}(s)$ and $\mathit{else}(s)$ to refer to the MDP state lifted from the respective ADD successor node.
We define $AP = \llbracket\tau \rrbracket \cup \{A, R\}$, and we define the labeling function $L$ as follows:
$$L(s) = \begin{cases}
	\{v, A\} &, \text{$s$ lifted from a terminal node with value $\mathit{val}(s) = v \neq R$}\\
	\{R\} &, \text{$s$ lifted from a terminal node with value $\mathit{val}(s) = R$}\\
	\emptyset & \text{otherwise}
\end{cases}$$
Here, states which were lifted from terminal nodes for values $v\in \llbracket\tau \rrbracket$ are labeled with their corresponding output in the ADD and the proposition $A$ to model an accepted run.
States lifted from terminal nodes containing the value $R$ are labeled with the proposition $R$.

Now consider some \NDICE expression $e$ with $\{\} \vdash e \colon \tau \leadsto (\dot{\varphi}, \gamma, t)$, as well as the unnormalized MDP $\mathcal{M}^u_e$ lifted from an ADD for $(\dot{\varphi}, \gamma, t)$.
One can prove that the \textit{unnormalized} semantics of $e$ are closely related to reachability probabilities in $\mathcal{M}^u_e$.
However, the central inference question for PPLs is to infer the behavior of the program under the condition that all observations succeeded, which corresponds to the \textit{normalized} semantics of $e$.
Hence, we are not interested in the likelihood to just reach a particular value $v$, but we are instead interested in the likelihood to reach a particular $v$ under the condition that all observations pass, as modeled by reaching any state labeled with $A$.

We are thus interested in the conditional reachability probability $Pr^{max}_{\mathcal{M}^u_e} (\lozenge v \;|\; \lozenge A)$.
As discussed at the end of \Cref{sec:prelims_MDP}, there are two approaches to determine these conditional probabilities:
The reduction to ordinary reachability properties by \citet{baier2014computing}, and the recent, bisection-based approach implemented by the Storm Model Checker \citeBisection.
We will be using the bisection-based algorithm, as it is highly suited for our specific use-case.
Namely, the algorithm is very efficient on acyclic MDPs, as our MDPs naturally are.
Furthermore, the runtime of the algorithm is agnostic to the probability of an accepting run, hence the runtime of the inference process will be unaffected by conditioning with rare events.
This is a desirable trait, since inference engines commonly struggle to infer probabilities which are conditioned by low-probability observations.


 
 \subsection{Soundness}
 It is now time to formally consider the connection between the semantics of a $\NDICE$ program $p$ and the unnormalized MDP lifted from its compiled Boolean formulas.
 In practice it is infeasible to infer all distributions in $\llbracket p \rrbracket_D$, and thus we provide an upper bound for their behavior instead.
 Formally, we show that the maximal probability to terminate with a value $v$ among all distributions in $\llbracket p \rrbracket_D$ is equal to the maximal conditional probability to reach a state labeled with $v$ in the unnormalized MDP lifted from the compiled formulas of $e$:

\begin{theorem}[Inference Soundness]\label{ProgramSoundness}
	Let $p$ be a \NDICE program  
	%with free variables $x_1,\dots,x_n$ 
	and let $\{\}, \{\} \vdash p \colon \tau \leadsto (\dot{\varphi}, \gamma, t)$. 
	Then the reduced ADD for $(\dot{\varphi}, \gamma, t)$ can be lifted to an unnormalized MDP $\mathcal{M}_p$ so that:
	$$ Pr^{max}_{M_p} (\lozenge v \,|\, \lozenge A) = \max \left\{d(v) \,|\, d \in \llbracket p\rrbracket^{\{\}}_D\right\}$$
\end{theorem}
\begin{proof}
	 \ext
	 {See \Cref{ProgramAppendix}, where the theorem is restated as \Cref{app:ProgramSoundness}.}
	 {We refer to the extended version \citeExtended.}
\end{proof}
%
%
\noindent%
%
%
The theorem above refers to the \textit{reduced} ADD for the formulas, meaning we may use the smallest ADD.
Still, the construction of this ADD is PSPACE-hard \cite{holtzen2020scaling}, and the final ADD itself may be exponentially large in the number of flip operations in the expression.
Additionally, we then still need to lift the ADD to an MDP and determine the necessary reachability probabilities.
Lifting the ADD to an MDP can be done in a single pass, while determining conditional reachability probabilities takes polynomial time relative to the MDP's size \cite{puterman2014markov, baier2014computing}.

On paper, this makes \NDICE inference very expensive.
However, \citet{holtzen2020scaling} show how the \DICE compilation naturally exploits structural properties like conditional independence, determinism or context-specific independence to build compact BDDs quickly.
\NDICE reuses the compilation of \DICE, and thus \NDICE exploits the same program structures for compact ADDs.
While we refer to \citet{holtzen2020scaling} for details on structures exploited by \DICE, we still want to briefly review one central source of efficiency for \DICE and \NDICE: Conditional independence. 


Intuitively, variables $x$ and $y$ are conditionally independent given $z$, if the behavior of $x$ does not depend on $y$ if a value for $z$ has been fixed.
Programs contain conditional independence naturally since it arises from constructs like \texttt{let} statements or function calls, as seen in the introductory program from \Cref{fig:plane}:
Here, all calls to \texttt{step} are conditionally independent of each other given their return values.
After all, subsequent function calls and the behavior of the program only depend on the functions output.
Hence, all intermediate values, like e.g. the flip operations in the \texttt{move} function, are conditionally independent of each other given the output of \texttt{step}.
This conditional independence prevents an exponential blow-up during the ADD construction, as each function call only adds another conditionally independent \enquote{layer} of nodes.



